% BHU PYQ -- Manually transcribed & error-corrected
% Paper: BCF-114 : Business Organisation
% Exam: B.Com. (Hons.) FMM Semester I Examination, 2016-17

\documentclass[11pt,a4paper]{article}
\usepackage[a4paper,top=2.5cm,bottom=2.5cm,left=2.8cm,right=2.8cm,headheight=14pt]{geometry}
\usepackage{amsmath,amssymb}
\usepackage[T1]{fontenc}
\usepackage[utf8]{inputenc}
\usepackage{lmodern,microtype,fancyhdr,enumitem,hyperref,lastpage,parskip}

\hypersetup{
    colorlinks=true,
    urlcolor=black,
    linkcolor=black,
    pdftitle={BCF-114: Business Organisation -- BHU B.Com. (Hons.) FMM Sem I 2016-17}
}

\pagestyle{fancy}
\fancyhf{}
\renewcommand{\headrulewidth}{0.4pt}
\renewcommand{\footrulewidth}{0.4pt}
\fancyhead[L]{\small\textit{Banaras Hindu University}}
\fancyhead[R]{\small\textit{BCF-114: Business Organisation}}
\fancyfoot[C]{\small Page \thepage\ of \pageref{LastPage}}

\newlist{parts}{enumerate}{2}
\setlist[parts,1]{label=(\alph*),leftmargin=*,itemsep=5pt,topsep=3pt}
\setlist[parts,2]{label=(\roman*),leftmargin=*,itemsep=3pt,topsep=2pt}
\newcommand{\pts}[1]{\hfill\textit{[#1]}}

\begin{document}

\begin{center}
  {\large\bfseries BANARAS HINDU UNIVERSITY}\\[2pt]
  {\normalsize B.Com. (Hons.) FMM --- Semester I Examination, 2016-17}\\[6pt]
  \rule{0.8\linewidth}{0.5pt}\\[6pt]
  {\large\bfseries Paper No. BCF-114: Business Organisation}\\[4pt]
  {\normalsize Time: 3 Hours\hfill Maximum Marks: 70}
\end{center}

\rule{\linewidth}{0.4pt}

\smallskip
\noindent\textit{Instructions: Answer all questions. All questions carry equal marks (14 marks each).}

\bigskip

\noindent\textbf{Question 1.}
What do you mean by Industrial Organization? Discuss its scope clearly.\pts{14}

\centerline{\textbf{OR}}

Define entrepreneurship and discuss its role in economic development.\pts{14}

\medskip\rule{\linewidth}{0.2pt}

\noindent\textbf{Question 2.}
Describe the factors which affect the location of an industrial unit.\pts{14}

\centerline{\textbf{OR}}

Define plant layout and briefly explain its main principles.\pts{14}

\medskip\rule{\linewidth}{0.2pt}

\noindent\textbf{Question 3.}
``The optimum size of a firm is determined through the reconciliation of different optima.'' Explain this process of reconciliation.\pts{14}

\centerline{\textbf{OR}}

Describe the causes responsible for industrial combination in the present era.\pts{14}

\medskip\rule{\linewidth}{0.2pt}

\noindent\textbf{Question 4.}
Define the terms merger and acquisition. What are the various types of business merger?\pts{14}

\centerline{\textbf{OR}}

``Sound organization is the base for the success of a business house.'' Elucidate this statement and explain the various types of business organization found in the present business world.\pts{14}

\medskip\rule{\linewidth}{0.2pt}

\noindent\textbf{Question 5.}
Discuss the various sources of long-term finance in India.\pts{14}

\centerline{\textbf{OR}}

Write short notes on any two of the following:\pts{14}
\begin{parts}
  \item Debenture and term loan
  \item Preference share capital
  \item Retained earnings
\end{parts}

\vfill
\rule{\linewidth}{0.4pt}
\begin{center}\small
\href{https://bhu-exam.vercel.app/}{Click here to visit our website}\\[4pt]
\href{https://me-aryan.vercel.app/}{Click here to visit our contributor}\\[4pt]
\href{https://quashan-me.vercel.app/}{Click here to visit our contributor}
\end{center}

\end{document}
